\subsection{谱隙驱动的时间尺度与读出稳定}\label{subsec:pom-spectral-gap}
\begin{corollary}[谱隙控制的读出时间尺度]\label{cor:pom-spectral-gap-time}
设 $M_K(u)$ 原始，主特征值为 $\lambda_1(u)$，其余特征值最大模为 $\lambda_2(u)$。由定理 \ref{thm:kernel-weighted-prime-orbit}，
\[
B_n(u)=\frac{\lambda_1(u)^n}{n}+O\!\left(\frac{|\lambda_2(u)|^n}{n}\right),
\]
从而相对误差满足
\[
\left|\frac{nB_n(u)}{\lambda_1(u)^n}-1\right|
\le C\left(\frac{|\lambda_2(u)|}{\lambda_1(u)}\right)^n.
\]
给定目标精度 $\varepsilon\in(0,1)$，满足
\[
n\ \ge\ \frac{\log(C/\varepsilon)}{\log\!\bigl(\lambda_1(u)/|\lambda_2(u)|\bigr)}
\]
即可保证相对误差不超过 $\varepsilon$。若进一步满足核版 RH/GRH 条件（定义 \ref{def:kernel-rh}），则可用 $\lambda_1(u)^{1/2}$ 替代 $|\lambda_2(u)|$，从而得到平方根级衰减尺度（推论 \ref{cor:kernel-rh-sqrt-error}）。
\end{corollary}

\begin{remark}[自对偶与完成化偶性]\label{rem:pom-duality-even}
若同步核满足附录 \ref{app:kernel-self-dual} 的自对偶关系，则完成化后的 $\zeta$ 函数在 $\theta$ 变量上呈严格偶性。于是其对 $\theta$ 的奇阶导数在 $\theta=0$ 处消失，这对代价势的涨落统计施加对称性约束；相应的可调性仅能通过破坏该自对偶结构实现。
\end{remark}

\subsubsection{内生时间相关：谱隙给出的相关长度尺度（无外部输入）}

本小节把“谱隙”从 prime-orbit 误差控制进一步解释为\textbf{时间相关的衰减尺度}。这里强调“无外部输入”：考虑的随机性来自核自身的 Parry/平衡测度（即由有限维传输算子诱导的平稳马尔可夫结构），而非外加噪声或外部驱动。

\begin{proposition}[谱隙 \texorpdfstring{$\Rightarrow$}{=>} 时间相关指数衰减（可审计）]\label{prop:pom-time-corr-gap}
设 $B_\theta$ 为有限状态同步核的倾斜传输矩阵族（例如对输出位势 $e$ 的一维倾斜），并假设对给定 $\theta$ 矩阵 $B_\theta$ 原始。记
\(
\lambda_1(\theta)=\rho(B_\theta),\quad
\Lambda(\theta)=\max\{ |\lambda|:\ \lambda\in\mathrm{spec}(B_\theta),\ |\lambda|<\lambda_1(\theta)\}
\)，并令
\[
\rho(\theta):=\frac{\Lambda(\theta)}{\lambda_1(\theta)}\in[0,1),
\qquad
g(\theta):=-\log \rho(\theta)\in(0,+\infty).
\]
则对任意依赖有限窗口的可观测量 $F,G$（以块提升化为单时刻函数后），在 $B_\theta$ 的 Parry/平衡测度下存在常数 $C_{F,G}(\theta)<\infty$ 使得对所有 $n\ge 0$ 有
\[
\bigl|\mathrm{Cov}(F,\,G\circ\sigma^n)\bigr|\ \le\ C_{F,G}(\theta)\,\rho(\theta)^n.
\]
因此可把
\(
\tau(\theta):=1/g(\theta)
\)
解释为“相关时间”，并可把 $\tau_{1/2}(\theta):=\log 2\,/\,g(\theta)$ 解释为半衰期（指数包络衰减到 $1/2$ 的时间）。
\end{proposition}

\begin{proof}
对原始非负矩阵 $B_\theta$，Perron--Frobenius 理论给出唯一主特征值 $\lambda_1(\theta)>0$ 及相应左右正特征向量。由 Parry 测度的标准构造（Doob 变换）可把协方差写为
\(
\mathrm{Cov}(F,G\circ\sigma^n)=\langle f,\,\mathcal{T}_\theta^n g\rangle
\)
的形式，其中 $\mathcal{T}_\theta$ 为归一化后的转移算子（谱半径 $1$），且其非平凡谱半径为 $\rho(\theta)$。对中心化函数，主特征模态消失，仅剩余谱贡献，从而得到指数衰减界。全部量均为有限维线性代数对象，可直接由矩阵谱与向量内积审计。
\end{proof}

\input{sections/generated/fig_sync_kernel_time_correlation}

\subsection{新结论簇 IV：指数尺度饱和与常数级指纹}\label{subsec:pom-lambda-saturation}
\begin{proposition}[指数尺度的输入饱和]\label{prop:pom-lambda-saturation}
对确定性、完备读入的同步核，Perron 根 $\lambda$ 主要由输入骨架决定，而核方向仅贡献更小模的谱修正。具体地，40 状态真实输入核满足 $\rho(M)=\rho(A)=\varphi^2$（命题 \ref{prop:real-input-40-skeleton-factor}），10 状态同步核满足 $\lambda=3$（附录 \ref{app:sync-kernel-basic} 的行列式分解）。因此 $\log\lambda$ 更接近“输入分支熵”的饱和尺度，而非机制差异的细粒度指纹。
\end{proposition}

\begin{remark}[区分机制的三类指纹]\label{rem:pom-mechanism-fingerprint}
当主尺度 $\lambda$ 饱和时，真正区分运算机制的对象来自常数级与扭曲谱层：
\begin{itemize}
  \item \textbf{有限部曲线 $\log\mathfrak{M}(\theta)$：}在同一切片上与 $P(\theta)=\log\rho(M_\theta)$ 并列记录，用以刻画常数级漂移（见 \ref{subsec:chi-ldp} 中的对比注）。
  \item \textbf{Dirichlet--Mertens 常数张量：}对按同余类分层的 primitive 轨道计数给出常数级偏置（见 \ref{subsec:arity-dirichlet-mertens-tensor} 及相关附录）。
  \item \textbf{最坏扭曲比 $\rho/\lambda$：}由 Chebotarev 型误差控制序投影混合时间尺度（定理 \ref{thm:kernel-chebotarev-exp} 与推论 \ref{cor:pom-order-mixing}）。
\end{itemize}
这三类量与 $\lambda$ 正交互补，构成“运算机制”的可审计指纹；因此当 $\lambda$ 被输入分支锁定时，核之间的差异必然落在 $\log\mathfrak{M}$ 与 Dirichlet--Mertens 张量等有限部接口上。
\end{remark}

\subsection{真实输入 40 状态核的 RH-sharp 谱界与 toy RH}\label{subsec:pom-rh-sharp}
\begin{theorem}[RH-sharp/Ramanujan 谱界]\label{thm:pom-rh-sharp}
真实输入 40 状态核的行列式分解
\[
\det(I-zM)=(1-z)^2(1+z)^3(1-3z+z^2)(1+z-z^2)
\]
（见附录 \ref{app:real-input-40-finite-rh}）给出非零谱
\[
\mathrm{spec}(M)\setminus\{0\}
=\{\lambda,\ -\sqrt{\lambda},\ 1^{(2)},\ (-1)^{(3)},\ \lambda^{-1},\ \lambda^{-1/2}\},
\qquad \lambda=\varphi^2.
\]
因此除 $\lambda$ 外的全部非零特征值 $\nu$ 满足
\[
|\nu|\le \sqrt{\lambda}=\varphi,
\]
且仅有一个非平凡特征值达界：$\nu_2=-\sqrt{\lambda}$。
\end{theorem}

\begin{corollary}[平方根振荡的最优性]\label{cor:pom-rh-sharp-an}
设 $a_n=\Tr(M^n)$，则
\[
a_n=\lambda^n+(-\sqrt{\lambda})^n+O(1),
\]
并且误差项不能改进到 $o(\lambda^{n/2})$。等价地，动力学版的“Chebyshev 平方根误差”在本核上达到临界最优。
\end{corollary}

\begin{proposition}[\texorpdfstring{$\zeta$}{zeta} 的结构性分解：输入乘积核 \texorpdfstring{$\times$}{x} \((-z)\) 扭曲黄金因子 \texorpdfstring{$\times$}{x} 平凡周期因子]\label{prop:pom-zeta-factorization-40}
令黄金均值一阶 Markov 骨架矩阵为
$$
A:=
\begin{pmatrix}
1&1\\
1&0
\end{pmatrix},
$$
并令 \(M\) 为真实输入 40 状态核的邻接矩阵（附录 \ref{app:real-input-40-abel}）。
则有精确代数分解
$$
\det(I-zM)
=(1-z)^2(1+z)\cdot \det\!\bigl(I-z(A\otimes A)\bigr)\cdot \det(I+zA).
$$
等价地，
$$
\zeta_M(z)
=\zeta_{A\otimes A}(z)\cdot \zeta_A(-z)\cdot \frac{1}{(1-z)^2(1+z)},
$$
其中 \(\zeta_B(z):=1/\det(I-zB)\) 为有限型子移位的 \(\zeta\)-函数。
\end{proposition}
\begin{proof}
矩阵 \(A\) 的特征值为 \(\varphi\) 与 \(-\varphi^{-1}\)，故 \(A\otimes A\) 的特征值为
\(\{\varphi^2,-1,-1,\varphi^{-2}\}\)，从而
$$
\det\!\bigl(I-z(A\otimes A)\bigr)=(1-\varphi^2 z)(1+z)^2(1-\varphi^{-2}z)=(1+z)^2(1-3z+z^2).
$$
同时
\(
\det(I+zA)=1+z-z^2
\)。
将以上两式与 \((1-z)^2(1+z)\) 相乘，即得
$$
(1-z)^2(1+z)^3(1-3z+z^2)(1+z-z^2),
$$
与定理 \ref{thm:pom-rh-sharp} 的 \(\det(I-zM)\) 因式分解一致。证毕。
\end{proof}

\begin{remark}[隐藏的 \texorpdfstring{$(-A)$}{(-A)} 因子：toy RH 临界线晶格的结构来源]\label{rem:pom-minusA-factor-critical-line}
命题 \ref{prop:pom-zeta-factorization-40} 的第三因子可进一步改写为
$$
\det(I+zA)=\det\!\bigl(I-z(-A)\bigr),
$$
从而
$$
\boxed{\;
\det(I-zM)
=(1-z)^2(1+z)\cdot \det\!\bigl(I-z(A\otimes A)\bigr)\cdot \det\!\bigl(I-z(-A)\bigr).
\;}
$$
等价地，在 \(\zeta\)-函数层面，
$$
\boxed{\;
\zeta_M(z)
=\frac{1}{(1-z)^2(1+z)}\cdot \zeta_{A\otimes A}(z)\cdot \zeta_{-A}(z),
\qquad
\zeta_{-A}(z):=\frac{1}{\det(I-z(-A))}=\frac{1}{\det(I+zA)}.
\;}
$$
这一改写把 toy RH 的“临界线晶格”来源完全剥离出来：因为 \(-A\) 的 Perron 根为 \(-\varphi\)，而主尺度为 \(\lambda=\varphi^2\)，故唯一达界特征值
\(
-\sqrt{\lambda}=-\varphi
\)
并非来自 \(40\) 状态复杂结构的偶然共振，而是由\textbf{黄金均值邻接算子 \(A\) 的符号翻转因子 \((-A)\)}必然嵌入所致。
因此开带 \(0<\Re(s)<1\) 中的非平凡极点列（定理 \ref{thm:pom-toy-rh}）可视为 \(\det(I-z(-A))=0\) 在 Dirichlet 坐标 \(z=\lambda^{-s}\) 下的直接像。
\end{remark}

\begin{remark}[置换 lift 观点：Ramanujan/RH-sharp 作为“新本征值界”]\label{rem:pom-permutation-lift-ramanujan}
附录 \ref{prop:real-input-40-skeleton-factor} 给出：真实输入 \(40\) 状态核可写成 \(4\times 4\) 的块矩阵，非零块均为由确定性转移诱导的 \(10\times 10\) 置换矩阵 \(B^{(d)}\)，并且常纤维子空间
\(
U=\RR^{4}\otimes \mathrm{span}\{\mathbf{1}\}
\)
在 \(M\) 作用下不变且 \(M|_{U}\simeq A\otimes A\)。
这正是图谱语言中的\textbf{置换 lift（覆盖）}：基图由 \(A\otimes A\) 给出，每条基边携带一个置换（本核由 \(d\in\{0,1,2\}\) 诱导），从而在纤维上装配得到提升图的邻接算子 \(M\)（参见 \cite{StarkTerras1996GraphCoverings,TerrasZetaGraphs} 及图 lift 的谱理论综述）。
\par
在该 lift 观点下，\(M\) 的谱自然分为两类：
\begin{itemize}
  \item \textbf{旧谱（base/水平谱）：}来自 \(U\) 上的作用，即基图 \(A\otimes A\) 的谱；
  \item \textbf{新谱（vertical/纤维正交谱）：}来自 \(U^\perp=\RR^{4}\otimes \mathbf{1}^\perp\) 的作用，其全部内容由“纵向残差因子”承载（附录 \ref{prop:real-input-40-skeleton-factor} 的 \(\det_{\mathrm{vert}}\)）。
\end{itemize}
因此定理 \ref{thm:pom-rh-sharp} 的 Ramanujan/RH-sharp 半径界
\(
\max_{\mu\neq\lambda}|\mu|\le \sqrt{\lambda}
\)
可以被理解为对 lift 的\textbf{新本征值}的统一压制界：把全部 vertical 模态的谱半径压在 \(\sqrt{\lambda}\) 以内，并且本核达到饱和仅由隐藏因子 \((-A)\) 强制（注 \ref{rem:pom-minusA-factor-critical-line}）。
这给出一个面向“核版 RH/GRH”的结构管道：把 RH/GRH 的谱界问题转译为 lift/覆盖的谱工程（构造/排除 near-coboundary 扭曲，使新谱受控），从而与 Ramanujan 图与 lift 的新谱控制问题同型（见 \cite{BiluLinial2006LiftsDiscrepancy,Friedman2008AlonSecondEigenvalue,LinialPuder2010WordMapsLifts} 的代表性结果）。
\end{remark}

\begin{corollary}[有限部常数的可加分解：输入 \texorpdfstring{$\oplus$}{+} 扭曲 \texorpdfstring{$\oplus$}{+} 平凡]\label{cor:pom-logM-split-40}
沿命题 \ref{prop:pom-zeta-factorization-40} 的 \(\zeta\)-分解
\(
\zeta_M=\zeta_{\mathrm{in}}\zeta_{\mathrm{twist}}\zeta_{\mathrm{triv}}
\)，其中主极点 \(z_\star=1/\lambda=\varphi^{-2}\) 只来自 \(\zeta_{\mathrm{in}}\)。
于是 Abel/Hadamard 有限部常数满足\textbf{精确可加分解}
\[
\log\mathfrak{M}_{\mathrm{tot}}
\;=\;
\log\mathfrak{M}_{\mathrm{in}}
\;+\;
\log\mathfrak{M}_{\mathrm{twist}}
\;+\;
\log\mathfrak{M}_{\mathrm{triv}},
\]
其中 \(\log\mathfrak{M}_{\mathrm{in}}\) 为 \(\zeta_{\mathrm{in}}\) 的有限部常数，而对解析因子 \(f\in\{\zeta_{\mathrm{twist}},\zeta_{\mathrm{triv}}\}\) 定义其“有限部贡献”为
\[
\log\mathfrak{M}_f
:=
\log f(z_\star)
\;+\;
\sum_{k\ge 2}\frac{\mu(k)}{k}\log f(z_\star^k).
\]
相应的数值拆分见表 \ref{tab:real_input_40_finite_part_split}（脚本 \texttt{scripts/exp\_real\_input\_40\_finite\_part\_split.py} 可复现）。
\end{corollary}

\begin{proof}
由定理 \ref{thm:abel-mertens-universal}，对任意具有简单主极点 \(z_\star\) 的 \(\zeta\)-函数有
\(
\log\mathfrak{M}=\log C+\sum_{k\ge 2}\mu(k)\log\zeta(z_\star^k)/k
\)。
这里 \(\zeta_M=\zeta_{\mathrm{in}}\zeta_{\mathrm{twist}}\zeta_{\mathrm{triv}}\)，且 \(\zeta_{\mathrm{twist}},\zeta_{\mathrm{triv}}\) 在 \(z=z_\star\) 解析非零，因此主留数满足
\(
C_{\mathrm{tot}}=C_{\mathrm{in}}\zeta_{\mathrm{twist}}(z_\star)\zeta_{\mathrm{triv}}(z_\star)
\)。
再将 \(\log\zeta_M=\log\zeta_{\mathrm{in}}+\log\zeta_{\mathrm{twist}}+\log\zeta_{\mathrm{triv}}\) 代入 Möbius 尾并按因子分组，即得可加分解式；对解析因子 \(f\) 的“有限部贡献”恰为
\(
\log f(z_\star)+\sum_{k\ge2}\mu(k)\log f(z_\star^k)/k
\)。
证毕。
\end{proof}

\begin{remark}[主极点与留数的代数闭式]\label{rem:pom-logM-split-40-exact}
在本核上有完全闭式：
\[
z_\star=\lambda^{-1}=\varphi^{-2}=\frac{3-\sqrt5}{2},
\qquad
C_{\mathrm{in}}=\frac{3}{10}+\frac{7\sqrt5}{50},
\qquad
\zeta_{\mathrm{triv}}(z_\star)=1+\frac{2\sqrt5}{5},
\qquad
\zeta_{\mathrm{twist}}(z_\star)=\frac{1+\sqrt5}{4},
\]
从而
\(
C_{\mathrm{tot}}=\frac{47+21\sqrt5}{100}
\)。
这些闭式与 \(\log\mathfrak{M}\) 的拆分常数一起由表 \ref{tab:real_input_40_finite_part_split} 给出可核验数值。
\end{remark}

\input{sections/generated/tab_real_input_40_finite_part_split}

\begin{proposition}[Lucas--平方闭式]\label{prop:pom-lucas-square}
令 Lucas 数
\(
L_n:=\varphi^n+(-\varphi^{-1})^n
\)。
则
\[
a_n=L_n^2+(-1)^n(L_n+1)+2.
\]
\end{proposition}
\begin{proof}
由附录 \ref{app:real-input-40-abel} 的闭式
\(
a_n=L_{2n}+(-1)^nL_n+2+3(-1)^n
\)
与恒等式 $L_{2n}=L_n^2-2(-1)^n$ 直接化简得到。
\end{proof}

\begin{proposition}[奇素长 primitive 的 Lucas 闭式]\label{prop:pom-primitive-prime-lucas}
令 \(p\) 为奇素数，\(L_n\) 为 Lucas 数。则真实输入 40 状态核的 primitive 轨道数满足
$$
p_p=\frac{L_p(L_p-1)}{p}.
$$
\end{proposition}
\begin{proof}
由 Möbius 反演（定义 \ref{def:pom-proj-prim}）对素数 \(p\) 有
\(
p_p=\frac{a_p-a_1}{p}
\)。
由命题 \ref{prop:pom-lucas-square} 且 \(p\) 为奇数（故 \((-1)^p=-1\)）得
\(
a_p=L_p^2-L_p+1
\)，并且 \(a_1=1\)。
代入即得 \(p_p=\frac{L_p^2-L_p}{p}\)。证毕。
\end{proof}

\begin{lemma}[Lucas 素数同余：\texorpdfstring{$L_p\equiv 1\pmod p$}{Lp=1 mod p}]\label{lem:pom-lucas-modp}
对任意素数 \(p\)，Lucas 数满足 \(L_p\equiv 1\ (\mathrm{mod}\ p)\)。
\end{lemma}
\begin{proof}
设 \(p\neq 5\)（当 \(p=5\) 直接计算 \(L_5=11\equiv 1\pmod 5\)）。在环 \(\mathbb{Z}[\sqrt5]\) 中，
由二项式展开与 \(\binom{p}{k}\equiv 0\pmod p\)（\(1\le k\le p-1\)）得
$$
(1+\sqrt5)^p+(1-\sqrt5)^p\equiv 2+(\sqrt5)^p+(-\sqrt5)^p\pmod p.
$$
当 \(p\) 为奇素数时右端后两项相消，故
\(
(1+\sqrt5)^p+(1-\sqrt5)^p\equiv 2\pmod p
\)；
当 \(p=2\) 直接验算亦成立。又由定义 \(2^p L_p=(1+\sqrt5)^p+(1-\sqrt5)^p\) 与费马小定理 \(2^{p-1}\equiv 1\pmod p\) 得
\(
2^p\equiv 2\pmod p
\)，因此 \(L_p\equiv 1\pmod p\)。证毕。
\end{proof}

\begin{remark}[整数性与 Wieferich 型稀有事件]\label{rem:pom-lucas-wieferich}
由命题 \ref{prop:pom-primitive-prime-lucas} 与引理 \ref{lem:pom-lucas-modp} 可见 \(L_p-1\) 总被 \(p\) 整除，从而 \(p_p\in\mathbb{Z}\) 的整数性在素长层自动成立。
更强的条件 \(p\mid p_p\) 等价于 \(p^2\mid(L_p-1)\)，对应 Lucas 序列的 Wieferich 型现象；这给出一个把“素长 primitive 退化”与经典数论稀有事件对接的接口。
\end{remark}

\begin{corollary}[primitive 奇偶分裂]\label{cor:pom-prim-odd-even}
记 $p_n$ 为 primitive 轨道数，则
\[
p_n=\frac{\lambda^n}{n}+\frac{\bigl((-1)^n-\mathbf{1}_{2\mid n}\bigr)\lambda^{n/2}}{n}
+O\!\left(\frac{\lambda^{n/3}}{n}\right).
\]
其中偶数情形的 $\lambda^{n/2}$ 级项在 Möbius 反演的 $d=1,2$ 两项间精确抵消（见附录 \ref{app:real-input-40-finite-rh}）。
\end{corollary}

\begin{remark}[偶长轨道的更低阶振荡]
写 $n=2m$，利用 $a_{m}=\lambda^{m}+(-\sqrt{\lambda})^{m}+O(1)$ 可进一步得到
\[
p_{2m}=\frac{\lambda^{2m}}{2m}-\frac{(-\sqrt{\lambda})^{m}}{2m}
+O\!\left(\frac{\lambda^{2m/3}}{m}\right),
\]
表明偶长 primitive 的首个显式次项规模为 $\lambda^{n/4}$。
\end{remark}

\begin{theorem}[Dirichlet 极点晶格与 toy RH]\label{thm:pom-toy-rh}
取 Dirichlet 变量 $z=\lambda^{-s}$ 与 $r=\lambda z=\lambda^{1-s}$。则 $\widehat{\zeta}(s)=\zeta_M(\lambda^{-s})$ 的全部极点由
\[
r\in\{1,\ \pm\lambda,\ \lambda^2,\ \varphi^3,\ -\varphi\}
\]
给出，并满足
\[
\Re(s)\in\left\{1,\tfrac12,0,-\tfrac12,-1\right\}.
\]
其中 $r=-\varphi$ 产生一条临界线晶格
\(
s=\tfrac12+\frac{(2k+1)\pi i}{\log\lambda}
\),
是开带 $0<\Re(s)<1$ 内的唯一极点列（见附录 \ref{app:real-input-40-finite-rh}）。
\end{theorem}

\begin{proposition}[输出压力的二、三阶累积量闭式]\label{prop:pom-pressure-cumulants}
令 $u=e^{\theta}$，$\lambda(u)=\rho(M(u))$，$P(\theta)=\log\lambda(e^\theta)$，并以
\(
G(\lambda,u)=0
\)
刻画 $\lambda(u)$（见附录 \ref{app:real-input-40-zeta-u}）。则
\[
P''(0)=\frac{1791}{200}-\frac{3983}{1000}\sqrt5,
\qquad
P'''(0)=\frac{1416369}{5000}-\frac{126691}{1000}\sqrt5,
\]
从而 $P''(0)\approx 0.0487412456$、$P'''(0)\approx -0.0158881374$。
\end{proposition}
\begin{proof}
对 $G(\lambda(u),u)=0$ 做隐函数求导，并用 $u=e^\theta$ 的链式法则把 $u$-导数转为 $\theta$-导数；代入 $u=1,\lambda=\varphi^2$ 即得。
\end{proof}

\begin{corollary}[零温极限常数]\label{cor:pom-pressure-zero-temp}
由附录 \ref{app:real-input-40-zeta-u} 可得 $\theta\to+\infty$ 时
\[
P(\theta)=\frac{\theta}{2}+\log c+o(1),
\]
其中 $c=\sqrt{y_{\max}}$，$y_{\max}$ 为方程 $y^4-6y^3+9y^2-y-1=0$ 的最大正根。
\end{corollary}

\subsection{素数影子到素数寄存器的桥接}\label{subsec:pom-prime-bridge}
\paragraph{有限轴束的扭曲与分解}
对有限群标签 $G$ 的扭曲矩阵族 $M_{K,\rho}(u)$ 与扭曲 $\zeta/L$-函数，命题 \ref{prop:kernel-artin-factorization} 给出 Artin 型因子化。该分解表明：有限轴束本质上是以 $\widehat G$ 为谱坐标系的“可审计纤维化”。

\paragraph{新增模轴的不可约性}
命题 \ref{prop:arity-335-n2-nondegenerate} 证明了第三轴 $N_2(\gamma)\bmod 5$ 无法由前两轴消去，因此新增模轴引入不可约的新自由度。这一事实把“增加轴束不是重标记”在核层面固定为可审计结论。

\begin{definition}[素数寄存器纤维的极限口径]\label{def:pom-profinite}
取模系统 $\{\ZZ/m\ZZ\}_{m\ge 1}$ 的射影极限
\[
\widehat{\ZZ}:=\varprojlim_m \ZZ/m\ZZ\ \cong\ \prod_{p\in\mathbb{P}}\ZZ_p,
\]
并将其作为“无限素数轴束”的极限对象。该极限把所有有限模轴束统一为同一纤维坐标体系，从而为“素数寄存器”提供结构性定义（参见 \cite{Apostol1976} 中关于 CRT 与 p-进分解的标准事实）。
\end{definition}

\begin{remark}[与 Dirichlet 角色的真实口径对齐：\texorpdfstring{$\widehat{\ZZ}^\times$}{Zhat^x}]\label{rem:pom-profinite-units}
在经典解析数论中，Dirichlet 角色作用在单位群 $(\ZZ/m\ZZ)^\times$ 上，其自然极限对象是
\[
\widehat{\ZZ}^{\,\times}:=\varprojlim_m (\ZZ/m\ZZ)^\times\ \cong\ \prod_{p\in\mathbb{P}}\ZZ_p^\times.
\]
本文在“同余轴束”层面以 $\widehat{\ZZ}$ 作为统一坐标，是为了把所有有限模 residue 读出纳入同一纤维化框架；当需要与 Dirichlet/Artin 的角色语言严格对齐时，只需在该框架内把读出限制到单位部分（并相应把角色基换成 $\widehat{\ZZ}^{\,\times}$ 的对偶坐标），即可得到与真实 GRH 更贴近的“全角色族”口径。
\end{remark}

\begin{remark}[从有限扭曲到无限寄存器]
有限群扭曲提供“有限模轴束”的可审计谱分解；而新增轴的不可约性意味着该分解在模数增大时严格增加自由度。因而 $\widehat{\ZZ}$ 的极限口径与素数寄存器纤维 $\mathcal{P}$ 在结构上是兼容的：前者是“模轴束的完备化”，后者是“乘法坐标的线性化轴束”。
\end{remark}

\begin{remark}[同余轴的逆极限与角色对角化]\label{rem:pom-congruence-invlimit}
将有限同余轴 $\{\ZZ/m\ZZ\}_{m\ge 1}$ 视作投影系统，则其逆极限给出
\(
\widehat{\ZZ}\cong\prod_{p}\ZZ_p
\)；
对应的角色基
\(
\chi_{k/m}(n)=e^{2\pi i kn/m}
\)
是 Pontryagin 对偶上的自然坐标。由此“有限同余轴 + 相位对角化 + 扭曲 $\zeta$”形成一条可审计的构造路径，极限即为“无限素数寄存器纤维”的严格实现。
\end{remark}

\begin{proposition}[profinite 角色对角化：Hilbert 空间上的 Artin 分解与核版 GRH 的算子化]\label{prop:pom-profinite-artin-hilbert}
令 \(K\) 为有限核（有向图/Markov 核），顶点集为 \(V\)，并令 \(\kappa:E(K)\to\ZZ\) 为一步可加势（例如 carry/碰撞/输出位势）。
取 profinite 同余轴 \(G_\infty:=\widehat{\ZZ}\)（或其单位部分 \(\widehat{\ZZ}^{\,\times}\)，见注 \ref{rem:pom-profinite-units}），并令
\(
g:E(K)\to G_\infty
\)
为一步同余 cocycle（即把每条边标记为某个 residue 增量；对有限模 \(m\) 的读出即为合成 \(G_\infty\to\ZZ/m\ZZ\)）。
在 Hilbert 空间
\[
\mathcal{H}:=\CC^{V}\otimes L^2(G_\infty)
\]
（\(G_\infty\) 上取 Haar 测度）上定义“群扩张邻接算子”
\[
\bigl(\mathcal{M}_u f\bigr)(j,h)
:=\sum_{e:i\to j} u^{\kappa(e)}\, f\!\bigl(i,\ h\cdot g(e)\bigr),
\qquad u>0,
\]
其中 \(f\in\mathcal{H}\) 视为 \(V\times G_\infty\) 上的平方可积函数，且 \(h\mapsto h\cdot g(e)\) 为右乘平移。
则：
\begin{enumerate}
  \item \textbf{（Fourier 分解与块对角化）} 由于 \(G_\infty\) 为紧致交换群，存在角色正交基分解
  \[
  L^2(G_\infty)\cong \bigoplus_{\chi\in\widehat{G_\infty}}\CC\,\chi,
  \]
  且在该分解下 \(\mathcal{M}_u\) 以酉等价方式化为直和
  \[
  \mathcal{M}_u\ \cong\ \bigoplus_{\chi\in\widehat{G_\infty}} M_\chi(u),
  \]
  其中每个块 \(M_\chi(u)\) 是 \(\CC^{V}\) 上的有限维扭曲矩阵：
  \[
  (M_\chi(u))_{i\to j}:=\sum_{e:i\to j} u^{\kappa(e)}\,\chi\!\bigl(g(e)\bigr).
  \]
  \item \textbf{（最坏扭曲谱半径的算子读出）} 记平凡角色为 \(1\)，并令
  \(
  \mathcal{H}_0:=\CC^V\otimes\mathrm{span}\{1\}
  \)
  及 \(\mathcal{H}_\perp:=\mathcal{H}_0^\perp\)。
  则
  \[
  \rho\!\left(\mathcal{M}_u|_{\mathcal{H}_\perp}\right)
  =\sup_{\chi\neq 1}\rho\!\left(M_\chi(u)\right),
  \qquad
  \rho\!\left(\mathcal{M}_u|_{\mathcal{H}_0}\right)=\rho\!\left(M_1(u)\right)=:\lambda_1(u).
  \]
  \item \textbf{（核版 GRH 的算子化）} 若核版 GRH（定义 \ref{def:kernel-rh} 的同型条件）在该 profinite 角色族上成立，即
  \(
  \sup_{\chi\neq 1}\rho(M_\chi(u))\le \lambda_1(u)^{1/2}
  \),
  则等价地有算子谱半径不等式
  \[
  \rho\!\left(\mathcal{M}_u|_{\mathcal{H}_\perp}\right)\ \le\ \lambda_1(u)^{1/2}.
  \]
\end{enumerate}
\end{proposition}
\begin{proof}
对有限模 \(m\) 的轴束，命题 \ref{prop:kernel-artin-factorization} 给出有限群版本的行列式因子化，其本质是正则表示的 Fourier/角色分解。
对紧致交换群 \(G_\infty\)，Pontryagin 对偶 \(\widehat{G_\infty}\) 为离散交换群，角色在 \(L^2(G_\infty)\) 中形成正交基；右乘平移 \(T_{g(e)}:f(h)\mapsto f(h\cdot g(e))\) 在该基上对角化，特征值为 \(\chi(g(e))\)。
因此 \(\mathcal{M}_u\) 在 \(L^2(G_\infty)\) 的 Fourier 基下逐角色约化为有限维块 \(M_\chi(u)\)，从而得到直和分解。谱半径恒等式由直和算子的谱为各块谱的并集给出。证毕。
\end{proof}

\begin{remark}[对照：Ihara--Ramanujan 机制的同构版]\label{rem:pom-profinite-ihara-ramanujan}
命题 \ref{prop:pom-profinite-artin-hilbert} 的结构与图论 Ihara-\(\zeta\) 的“图版 RH \(\Leftrightarrow\) Ramanujan 谱界”完全同型：
角色分解把覆盖/轴束上的算子对角化为扭曲块，而“平方根界”就是对全部非平凡块的统一谱半径控制。
本文的核版 RH/GRH（定义 \ref{def:kernel-rh}）正是把这一机制移植到“primitive 轨道计数/显式公式误差”的可审计场景中（参见注 \ref{rem:kernel-graph-rh}）。
\end{remark}

\begin{remark}[素数长度作为并行结构探针]\label{rem:pom-prime-length-probe}
在张量装配的全局核上，迹序列满足 $a_n^{\mathrm{glob}}(u)=a_n^{\mathrm{flow}}(u)^m$，primitive 谱由 Witt/Möbius 反演给出（推论 \ref{cor:par-add-global-lift-endpoints}）。当 $n=\ell$ 为素数时，
\[
p_\ell^{\mathrm{glob}}(u)=\frac{1}{\ell}\Bigl(a_\ell^{\mathrm{flow}}(u)^m-a_1^{\mathrm{flow}}(u^\ell)^m\Bigr),
\]
从而素数长度的 primitive 项直接隔离了“不可约循环”与“幂重复循环”的差异。该差异对并行张量化极其敏感，可作为隐藏并行度与结构分解的无偏探针。
\end{remark}

\subsection{投影词正规形与读出层级塔}\label{subsec:pom-normal-form}
\begin{theorem}[半环闭包的一次投影正规形]\label{thm:pom-one-fold-normal-form}
令 $t$ 为由 $\oplus,\otimes,0_X,1_X$ 生成的任意项，变量取值于 $X_\infty$。则存在多项式
\(
P\in\NN[x_1,\dots,x_k]
\)
使得对任意 $c_1,\dots,c_k\in X_\infty$，
\[
t(c_1,\dots,c_k)=Z\!\left(P\bigl(\mathrm{Val}(c_1),\dots,\mathrm{Val}(c_k)\bigr)\right).
\]
因此在半环闭包内，规范投影 $P_Z$ 可在末端统一调用一次，而中间步骤可完全留在过程层。
\end{theorem}
\begin{proof}
由定理 \ref{thm:add-definitional} 与 \ref{thm:mul-definitional}，$\mathrm{Val}$ 给出 $(X_\infty,\oplus,\otimes)$ 与 $(\NN,+,\times)$ 的半环同构。任意项 $t$ 在 $\NN$ 上对应一个多项式 $P$；于是
\(
\mathrm{Val}(t(c_1,\dots,c_k))=P(\mathrm{Val}(c_1),\dots,\mathrm{Val}(c_k))
\)。
再由 $Z=\mathrm{Val}^{-1}$ 得结论。证毕。
\end{proof}

\begin{remark}[纲领：读出层级塔]\label{rem:pom-hyper-iteration}
定理 \ref{thm:pom-one-fold-normal-form} 表明，在半环闭包内“规范投影只需末端一次”。进一步可沿定理 \ref{thm:composition-two-layer} 的思路构造读出层级塔：第 $0$ 层以合成 $\circ$ 的一次迭代读出对应加法，第 $1$ 层以“加法迭代”对应乘法；按同样方式递归定义第 $k$ 层的“运算”为第 $k{-}1$ 层的迭代读出。该层级塔在过程层不引入新原语，而将复杂度转移为读出层级的提升。
\end{remark}
